% C1 -- MATCHED TWIN-PLANT EVALUATION HARNESS.
% Justifies attributing the measured difference to the residual path alone.
\begin{tikzpicture}[node distance=6mm and 9mm]
  \node[blk] (ref) {shared mission\\generator $q_d(t)$};

  % --- baseline twin -------------------------------------------------------
  \node[clas,above right=7mm and 12mm of ref] (pd1) {quaternion PD\\$K_q,K_\omega$};
  \node[plant,right=of pd1]                   (pl1) {MuJoCo twin A};
  \node[meas,right=of pl1]                    (m1)  {$\bar\theta_q^{\rm PD}$,\;$\bar\tau$,\;$f_{\rm sat}$};

  % --- augmented twin ------------------------------------------------------
  \node[clas,below right=7mm and 12mm of ref] (pd2) {\emph{identical} PD\\same $K_q,K_\omega$};
  \node[neuro,right=of pd2]                   (rsd) {CMG$+$PES residual\\$\beta,\eta,\chi$};
  \node[plant,right=of rsd]                   (pl2) {MuJoCo twin B};
  \node[meas,right=of pl2]                    (m2)  {$\bar\theta_q^{\rm CMG}$,\;$\bar\tau$,\;$f_{\rm sat}$};

  \node[blk,right=12mm of m1] (cmp) {matched comparison\\
        $\Delta=1-\bar\theta_q^{\rm CMG}/\bar\theta_q^{\rm PD}$};

  \draw[sig] (ref.east) -- ++(5mm,0) |- (pd1.west);
  \draw[sig] (ref.east) -- ++(5mm,0) |- (pd2.west);
  \draw[sig] (pd1) -- (pl1); \draw[sig] (pl1) -- (m1);
  \draw[sig] (pd2) -- (rsd); \draw[sig] (rsd) -- (pl2); \draw[sig] (pl2) -- (m2);
  \draw[sig] (m1) -- (cmp);
  \draw[sig] (m2.east) -- ++(6mm,0) |- (cmp.south);
  \draw[sigd] (pl1.north) -- ++(0,5mm) -| (pd1.north);
  \draw[sigd] (pl2.south) -- ++(0,-5mm) -| (pd2.south);

  \node[note,anchor=north west,text width=0.92\linewidth] at ($(pd2.south west)+(-2.6cm,-11mm)$)
    {Both twins share reference, duration, integration step, actuator scale,
     allocator, initial condition, safety limits and logging. The residual
     path is the only difference, so the reported improvement is
     attributable to it rather than to re-tuning.};
\end{tikzpicture}
